% ==============================================================================
%  5.1  El modelo de regresión lineal en notación matricial
% ==============================================================================
\section{El modelo en notación matricial}
\label{sec:t5-modelo}

El modelo de regresión lineal simple (\cref{def:rls}) se formula en forma escalar como
\[
  Y_i = \beta_0+\beta_1X_i+\varepsilon_i, \qquad \varepsilon_i\iid\Normal(0,\sigma^2), \quad i=1,\ldots,n,
\]
es decir, como un sistema de $n$ ecuaciones, una por observación:
\[
  \begin{aligned}
    Y_1 &= \beta_0+\beta_1X_1+\varepsilon_1, \\
    Y_2 &= \beta_0+\beta_1X_2+\varepsilon_2, \\
        &\;\;\vdots \\
    Y_n &= \beta_0+\beta_1X_n+\varepsilon_n .
  \end{aligned}
\]
En \textbf{forma matricial},
\[
  \underbrace{\begin{pmatrix} Y_1\\ Y_2\\ \vdots\\ Y_n \end{pmatrix}}_{\vect{Y}_{n\times1}}
  =
  \underbrace{\begin{pmatrix} 1 & X_1\\ 1 & X_2\\ \vdots & \vdots\\ 1 & X_n \end{pmatrix}}_{\vect{X}_{n\times2}}
  \underbrace{\begin{pmatrix} \beta_0\\ \beta_1 \end{pmatrix}}_{\vect{\beta}_{2\times1}}
  +
  \underbrace{\begin{pmatrix} \varepsilon_1\\ \varepsilon_2\\ \vdots\\ \varepsilon_n \end{pmatrix}}_{\vect{\varepsilon}_{n\times1}},
\]
o, abreviadamente, $\vect{Y}=\vect{X}\vect{\beta}+\vect{\varepsilon}$.

\begin{definicion}[Modelo de regresión lineal en notación matricial][def:t5-modelo]
  \[
    \vect{Y} = \vect{X}\vect{\beta}+\vect{\varepsilon},
    \qquad \vect{\varepsilon}\sim\Normal(\vect{0},\ \sigma^2\vect{I}_n),
  \]
  donde
  \begin{itemize}
    \item $\vect{Y}_{n\times1}$ es el \textbf{vector de respuestas};
    \item $\vect{\beta}_{2\times1}$ es el \textbf{vector de parámetros};
    \item $\vect{X}_{n\times2}$ es la \textbf{matriz de diseño}, conocida;
    \item $\vect{\varepsilon}_{n\times1}$ es el \textbf{vector de errores}, normal multivariante con
      media $\vect{0}$ y matriz de varianzas-covarianzas
      \[
        \Var(\vect{\varepsilon}) =
        \begin{pmatrix}
          \sigma^2 & 0 & \cdots & 0\\
          0 & \sigma^2 & \cdots & 0\\
          \vdots & \vdots & \ddots & \vdots\\
          0 & 0 & \cdots & \sigma^2
        \end{pmatrix}
        = \sigma^2\vect{I}_n .
      \]
  \end{itemize}
\end{definicion}

La forma diagonal de $\Var(\vect{\varepsilon})$ recoge a la vez la homogeneidad de varianzas
($\Var(\varepsilon_i)=\sigma^2$ para todo $i$) y la independencia de los errores
($\Cov(\varepsilon_i,\varepsilon_j)=0$ para $i\neq j$, que en el caso normal equivale a
independencia). Equivalentemente, $\vect{Y}\sim\Normal(\vect{X}\vect{\beta},\ \sigma^2\vect{I}_n)$,
que es la versión matricial de la \cref{prop:t2-dist-Y}.

El modelo es un caso particular del modelo lineal general $\vect{Y}=Z\vect{\theta}+\vect{U}$ del
Tema 1 (\cref{sec:t1-modelos}), con $Z=\vect{X}$, $\vect{\theta}=\vect{\beta}$ y
$\vect{U}=\vect{\varepsilon}$. Esta formulación permite extender directamente los resultados al
caso de varias variables explicativas, añadiendo columnas a $\vect{X}$ y componentes a
$\vect{\beta}$.
